\input{../tex-inputs/latex-leseansicht-vorspann}

\section[Theodor von Sosnosky an Arthur Schnitzler, 22. 6. 1900]{L04227 Theodor von Sosnosky an Arthur Schnitzler, 22. 6. 1900}
\nopagebreak\mylabel{L04227v}
\rehead{ }\normalsize\beginnumbering\briefempfaengerindex{Schnitzler, Arthur@\textsc{Schnitzler, Arthur}!zzzSosnosky, Theodor von@\emph{von Theodor von Sosnosky}!1900-06-221@{22. 6. 1900}|(be}
\toendnotes[C]{\smallbreak\pagebreak[2]}
\correspDesc{Versand  durch Theodor von Sosnosky am 22. 6. 1900 in Kremsmünster
\newline{}Erhalt  durch Arthur Schnitzler im Zeitraum [23. 6. 1900
                  – 27. 6. 1900?] in Wien}\toendnotes[C]{\smallbreak}
\Standort{DLA, A:Schnitzler, HS.1985.1.4640.}
\physDesc{Brief, 2 Blätter, 6 Seiten, 4645 Zeichen
\newline{}Handschrift: blaue Tinte, deutsche Kurrent}\toendnotes[C]{\smallbreak}
\pstart
           \raggedleft{}{\pb}dzt. \textsc{Kremsmünster\oindex{Kremsmünster@\textbf{Kremsmünster}, \emph{Verwaltungsgebiet}|pw}}, 22./6. 900.\pend
           
\pstart{}Sehr geehrter Herr Doktor!\pend\vspace{0.5em}
\pstart
           Soeben hab' ich \label{K_L04227-1v}\edtext{Ihre liebenswürdige
                  Sendung}{\lemma{\textnormal{\emph{Ihre … Sendung}}}\Cendnote{\textnormal{Es handelte sich um ein
                  Exemplar des \emph{Reigen}\pwindex{Schnitzler, Arthur 15.\,5.\,1862 Wien – 21.\,10.\,1931 ebd.@\textsc{Schnitzler, Arthur} (15.\,5.\,1862 Wien – 21.\,10.\,1931 ebd.), \emph{Schriftsteller, Mediziner}!Reigen. Zehn Dialoge@\strich\emph{Reigen. Zehn Dialoge}|pwk}, der im April 1900 in limitierter Charge als Privatdruck für ausgewählte Bekannte
                  publiziert worden war. Sosnosky\pwindex{Sosnosky, Theodor von 4.\,1.\,1866 Budapest – 3.\,2.\,1943 Wien@\textsc{Sosnosky, Theodor von} (4.\,1.\,1866 Budapest – 3.\,2.\,1943 Wien), \emph{Schriftsteller}|pwk} berichtete
                  nach Schnitzlers Tod über die Begebenheit:
                     »Ich bekam das Buch von einem persönlichen Freunde\pwindex{Schwarzkopf, Gustav 7.\,11.\,1853 Wien – 13.\,11.\,1939 ebd.@\textsc{Schwarzkopf, Gustav} (7.\,11.\,1853 Wien – 13.\,11.\,1939 ebd.), \emph{Schriftsteller}|pwv}{ }Schnitzlers, der auch mir befreundet war
                     und eines der zweihundert Exemplare\pwindex{Schnitzler, Arthur 15.\,5.\,1862 Wien – 21.\,10.\,1931 ebd.@\textsc{Schnitzler, Arthur} (15.\,5.\,1862 Wien – 21.\,10.\,1931 ebd.), \emph{Schriftsteller, Mediziner}!Reigen. Zehn Dialoge@\strich\emph{Reigen. Zehn Dialoge}|pwv} erhalten hatte, zur Lektüre und unterhielt mich köstlich.
                     Als Schnitzler dies erfuhr, erfreute er
                     mich durch die Zusendung eines Widmungsexemplars\pwindex{Schnitzler, Arthur 15.\,5.\,1862 Wien – 21.\,10.\,1931 ebd.@\textsc{Schnitzler, Arthur} (15.\,5.\,1862 Wien – 21.\,10.\,1931 ebd.), \emph{Schriftsteller, Mediziner}!Reigen. Zehn Dialoge@\strich\emph{Reigen. Zehn Dialoge}|pwv} des Buches.« (Theodor von Sosnosky\pwindex{Sosnosky, Theodor von 4.\,1.\,1866 Budapest – 3.\,2.\,1943 Wien@\textsc{Sosnosky, Theodor von} (4.\,1.\,1866 Budapest – 3.\,2.\,1943 Wien), \emph{Schriftsteller}|pwk}: \emph{Artur Schnitzler als Lyriker. Ein Gedenkblatt zu seinem 70.
                        Geburtstag}\pwindex{Sosnosky, Theodor von 4.\,1.\,1866 Budapest – 3.\,2.\,1943 Wien@\textsc{Sosnosky, Theodor von} (4.\,1.\,1866 Budapest – 3.\,2.\,1943 Wien), \emph{Schriftsteller}!Artur Schnitzler als Lyriker. Ein Gedenkblatt zu seinem 70. Geburtstag@\strich\emph{Artur Schnitzler als Lyriker. Ein Gedenkblatt zu seinem 70. Geburtstag}|pwk}. In: \emph{Neues Wiener
                        Journal}\pwindex{Neues Wiener Journal@\emph{Neues Wiener Journal}|pwk}, Nr. 13.825, 18. 5. 1932, S. 6–7). Der
                  vermittelnde Freund war Gustav
                  Schwarzkopf\pwindex{Schwarzkopf, Gustav 7.\,11.\,1853 Wien – 13.\,11.\,1939 ebd.@\textsc{Schwarzkopf, Gustav} (7.\,11.\,1853 Wien – 13.\,11.\,1939 ebd.), \emph{Schriftsteller}|pwk}.}}}\label{K_L04227-1} erhalten; ich beeile mich, Ihnen dafür meinen
               verbindlichſten Dank zu{ }ſagen und Ihnen die aufrichtige Verſicherung zu geben, daß es
               eine Überraſchung war, die mich im Gegenſatze zu den meiſten Überraſchungen – \uline{wirklich} erfreut hat.\pend
           
\pstart
           Wenn ich Ihnen hier ſage, daß mir das Buch\pwindex{Schnitzler, Arthur 15.\,5.\,1862 Wien – 21.\,10.\,1931 ebd.@\textsc{Schnitzler, Arthur} (15.\,5.\,1862 Wien – 21.\,10.\,1931 ebd.), \emph{Schriftsteller, Mediziner}!Reigen. Zehn Dialoge@\strich\emph{Reigen. Zehn Dialoge}|pwv}{ }ſehr gefallen hat, ſo \uline{könnte} das ja eine höfliche Erwiderung Ihrer Liebenswürdigkeit ſein; daß es
               aber die \uline{reine} Wahrheit iſt, kann Ihnen Herr \textsc{Schwarzkopf\pwindex{Schwarzkopf, Gustav 7.\,11.\,1853 Wien – 13.\,11.\,1939 ebd.@\textsc{Schwarzkopf, Gustav} (7.\,11.\,1853 Wien – 13.\,11.\,1939 ebd.), \emph{Schriftsteller}|pw}} beſtätigen; er hat es wohl ſchon gethan, und ich nehme an, daß mein ihm \label{K_L04227-2v}\edtext{\textsc{\begin{otherlanguage}{french}vis à vis\end{otherlanguage}}}{\lemma{\textnormal{\emph{vis à vis}}}\Cendnote{\textnormal{französisch: von Angesicht zu
                  Angesicht}}}\label{K_L04227-2} geäußertes unverhohlenes Entzücken über dieſe Skizzen\pwindex{Schnitzler, Arthur 15.\,5.\,1862 Wien – 21.\,10.\,1931 ebd.@\textsc{Schnitzler, Arthur} (15.\,5.\,1862 Wien – 21.\,10.\,1931 ebd.), \emph{Schriftsteller, Mediziner}!Reigen. Zehn Dialoge@\strich\emph{Reigen. Zehn Dialoge}|pwv} der Anlaß Ihrer ſo freundlichen
               Sendung {[}war{]}. Ich hatte Herrn \textsc{Schwarzkopf\pwindex{Schwarzkopf, Gustav 7.\,11.\,1853 Wien – 13.\,11.\,1939 ebd.@\textsc{Schwarzkopf, Gustav} (7.\,11.\,1853 Wien – 13.\,11.\,1939 ebd.), \emph{Schriftsteller}|pw}} ja erſucht, Sie von dem {\pb}großen Erfolg zu verſtändigen, den das
                  Buch\pwindex{Schnitzler, Arthur 15.\,5.\,1862 Wien – 21.\,10.\,1931 ebd.@\textsc{Schnitzler, Arthur} (15.\,5.\,1862 Wien – 21.\,10.\,1931 ebd.), \emph{Schriftsteller, Mediziner}!Reigen. Zehn Dialoge@\strich\emph{Reigen. Zehn Dialoge}|pwv} bei mir gehabt hat,
               denn ich hatte das Bedürfnis, Sie das wißen zu laßen, obſchon ich keineswegs so
               naiv-eitel war, zu glauben, daß mein Urteil für Sie von irgend einem Belang ſein
               könne. Wie mir eben einfällt, könnte dieſer Herrn \textsc{Schwrzkpf\pwindex{Schwarzkopf, Gustav 7.\,11.\,1853 Wien – 13.\,11.\,1939 ebd.@\textsc{Schwarzkopf, Gustav} (7.\,11.\,1853 Wien – 13.\,11.\,1939 ebd.), \emph{Schriftsteller}|pw}} gegebene Auftrag leicht als ein\strikeout{e} cachirter
               Wink, mir das Buch\pwindex{Schnitzler, Arthur 15.\,5.\,1862 Wien – 21.\,10.\,1931 ebd.@\textsc{Schnitzler, Arthur} (15.\,5.\,1862 Wien – 21.\,10.\,1931 ebd.), \emph{Schriftsteller, Mediziner}!Reigen. Zehn Dialoge@\strich\emph{Reigen. Zehn Dialoge}|pwv} zu ſenden,
               aufgefaßt werden; ich kann Ihnen aber mein \uline{Wort}
               geben, daß ich, als \introOben{}ich\introOben{} Herrn \textsc{Schw}.\pwindex{Schwarzkopf, Gustav 7.\,11.\,1853 Wien – 13.\,11.\,1939 ebd.@\textsc{Schwarzkopf, Gustav} (7.\,11.\,1853 Wien – 13.\,11.\,1939 ebd.), \emph{Schriftsteller}|pwv} bat, Ihnen meine Anſicht mitzuteilen, auch
               nicht den leiſesten derartigen Hintergedanken hatte. Dies geht auch daraus hervor,
               daß ich, als Herr \textsc{Schwkpf\pwindex{Schwarzkopf, Gustav 7.\,11.\,1853 Wien – 13.\,11.\,1939 ebd.@\textsc{Schwarzkopf, Gustav} (7.\,11.\,1853 Wien – 13.\,11.\,1939 ebd.), \emph{Schriftsteller}|pw}}  mir, (bevor er mir ſein Exemplar\pwindex{Schnitzler, Arthur 15.\,5.\,1862 Wien – 21.\,10.\,1931 ebd.@\textsc{Schnitzler, Arthur} (15.\,5.\,1862 Wien – 21.\,10.\,1931 ebd.), \emph{Schriftsteller, Mediziner}!Reigen. Zehn Dialoge@\strich\emph{Reigen. Zehn Dialoge}|pwv} lieh) angetragen, mir Ihr Buch\pwindex{Schnitzler, Arthur 15.\,5.\,1862 Wien – 21.\,10.\,1931 ebd.@\textsc{Schnitzler, Arthur} (15.\,5.\,1862 Wien – 21.\,10.\,1931 ebd.), \emph{Schriftsteller, Mediziner}!Reigen. Zehn Dialoge@\strich\emph{Reigen. Zehn Dialoge}|pwv} durch Sie ſelber zu verſchaffen, dies \introOben{}kurz\introOben{} ablehnte, da ich als ein für Sie Stockfremder keinen Anspruch darauf
               hätte.\pend
           
\pstart
           So viel zur Rechtfertigung!\pend
           
\pstart
           Über das Buch\pwindex{Schnitzler, Arthur 15.\,5.\,1862 Wien – 21.\,10.\,1931 ebd.@\textsc{Schnitzler, Arthur} (15.\,5.\,1862 Wien – 21.\,10.\,1931 ebd.), \emph{Schriftsteller, Mediziner}!Reigen. Zehn Dialoge@\strich\emph{Reigen. Zehn Dialoge}|pwv} bin ich der
               Anſicht, daß es Ihr \uline{bestes} Werk iſt – (ich kenne alle
              {\pb}Ihre gedruckten Sachen); ich bin nämlich \uline{nicht}
               der Anſicht, das ſeien graziöse pornographische Spielereien für
               Herren-Abende, ſondern ich{ }ſehe darin brillante Kabinetstücke erotiſcher Poeſie, die
               bleibenden großen Wert haben und es mehr verdienten, der Nachwelt erhalten zu bleiben
               als \textsc{Boccaccios\pwindex{Boccaccio, Giovanni 16.\,6.\,1313 Certaldo – 21.\,12.\,1375 ebd.@\textsc{Boccaccio, Giovanni} (16.\,6.\,1313 Certaldo – 21.\,12.\,1375 ebd.), \emph{Schriftsteller}|pw}} »\textsc{Decamerone\pwindex{Boccaccio, Giovanni 16.\,6.\,1313 Certaldo – 21.\,12.\,1375 ebd.@\textsc{Boccaccio, Giovanni} (16.\,6.\,1313 Certaldo – 21.\,12.\,1375 ebd.), \emph{Schriftsteller}!Decamerone@\strich\emph{Decamerone}|pw}}« oder das »\textsc{Heptameron\pwindex{de Navarre, Marguerite 11.\,4.\,1492 Angoulême – 21.\,12.\,1549 Odos@\textsc{de Navarre, Marguerite} (11.\,4.\,1492 Angoulême – 21.\,12.\,1549 Odos), \emph{Schriftstellerin}!Heptameron@\strich\emph{Heptameron}|pw}}« der Margarethe von Narvarra\pwindex{de Navarre, Marguerite 11.\,4.\,1492 Angoulême – 21.\,12.\,1549 Odos@\textsc{de Navarre, Marguerite} (11.\,4.\,1492 Angoulême – 21.\,12.\,1549 Odos), \emph{Schriftstellerin}|pw}. Wie aus dem
               Fortleben dieſer 2 Werke\pwindex{de Navarre, Marguerite 11.\,4.\,1492 Angoulême – 21.\,12.\,1549 Odos@\textsc{de Navarre, Marguerite} (11.\,4.\,1492 Angoulême – 21.\,12.\,1549 Odos), \emph{Schriftstellerin}!Heptameron@\strich\emph{Heptameron}|pwv}\pwindex{Boccaccio, Giovanni 16.\,6.\,1313 Certaldo – 21.\,12.\,1375 ebd.@\textsc{Boccaccio, Giovanni} (16.\,6.\,1313 Certaldo – 21.\,12.\,1375 ebd.), \emph{Schriftsteller}!Decamerone@\strich\emph{Decamerone}|pwv} hervorgeht, haben gerade erotische Bücher die Anwartschaft auf
               litterarische Langlebigkeit, das liegt in der erotischen Natur der Menschheit.\pend
           
\pstart
           Wenn nun die Erotik in ſo überaus ſinnreicher und graziöser Weiſe gehandhabt wird,
               wenn die Frivolität bei allem Übermut auf ſo tiefgründiger Seelenke{\geminationn}tnis beruht, dann hört jede Berechtigung für die
               freundliche Geringschätzung auf, die man ſonst erotischen Novelletten
               entgegenzubringen pflegt.\pend
           
\pstart
           Ich habe Ihr Buch\pwindex{Schnitzler, Arthur 15.\,5.\,1862 Wien – 21.\,10.\,1931 ebd.@\textsc{Schnitzler, Arthur} (15.\,5.\,1862 Wien – 21.\,10.\,1931 ebd.), \emph{Schriftsteller, Mediziner}!Reigen. Zehn Dialoge@\strich\emph{Reigen. Zehn Dialoge}|pwv} das erſte
              Mal unter den {\pb}günſtigſten Umständen gelesen: an demselben
                  Abend, an dem ich es erhalten, ſaß ich allein beim \label{K_L04227-3v}\edtext{Conzert\eventindex{Volksgartenrestaurant@\textbf{Volksgartenrestaurant}!?? [Konzert im Volksgartengebäude, das Theodor von Sosnosky besucht, als er gerade das Widmungsexemplar des Reigens erhalten hat]@?? [Konzert im Volksgartengebäude, das Theodor von Sosnosky besucht, als er gerade das Widmungsexemplar des Reigens erhalten hat]|pw} im Volksgarten-Gebäude\oindex{Wien@\textbf{Wien}!I., Innere Stadt@\textbf{I., Innere Stadt}!Volksgartenrestaurant@\textbf{Volksgartenrestaurant}, \emph{Veranstaltungsgebäude}|pw}}{\lemma{\textnormal{\emph{Conzert im Volksgarten-Gebäude}}}\Cendnote{\textnormal{Das Café-Restaurant\oindex{Wien@\textbf{Wien}!I., Innere Stadt@\textbf{I., Innere Stadt}!Volksgartenrestaurant@\textbf{Volksgartenrestaurant}, \emph{Veranstaltungsgebäude}|pwk} im Volksgarten\oindex{Wien@\textbf{Wien}!I., Innere Stadt@\textbf{I., Innere Stadt}!Volksgarten@\textbf{Volksgarten}, \emph{Park}|pwk}, das
                  seit 1899 von Johann Seidl\pwindex{Seidl, Johann @\textsc{Seidl, Johann}, \emph{Restaurantbesitzer}|pwk}
                  betrieben wurde, verfügte über einen Konzertsaal, in dem nahezu abendlich Programm
                  geboten wurde.}}}\label{K_L04227-3}. Um mir die Zeit zu vertreiben, nahm ich das Buch\pwindex{Schnitzler, Arthur 15.\,5.\,1862 Wien – 21.\,10.\,1931 ebd.@\textsc{Schnitzler, Arthur} (15.\,5.\,1862 Wien – 21.\,10.\,1931 ebd.), \emph{Schriftsteller, Mediziner}!Reigen. Zehn Dialoge@\strich\emph{Reigen. Zehn Dialoge}|pwv} vor und begann zu leſen, umklungen von
               hübſchen Melodien, umgeben von fröhlichen angeregten Menſchen, unter denen ſicher
               manches »ſüße Mädl\pwindex{Schnitzler, Arthur 15.\,5.\,1862 Wien – 21.\,10.\,1931 ebd.@\textsc{Schnitzler, Arthur} (15.\,5.\,1862 Wien – 21.\,10.\,1931 ebd.), \emph{Schriftsteller, Mediziner}!Reigen. Zehn Dialoge@\strich\emph{Reigen. Zehn Dialoge}|pwv}« mit ihrem
               »Freunde« und vielleicht auch manche Frau mit Gatten {\kaufmannsund}
               Hausfreund{ }ſaß. Es war das richtige Milieu, und ich konnte die Lektüre nicht laßen,
               eh’ ich das Buch\pwindex{Schnitzler, Arthur 15.\,5.\,1862 Wien – 21.\,10.\,1931 ebd.@\textsc{Schnitzler, Arthur} (15.\,5.\,1862 Wien – 21.\,10.\,1931 ebd.), \emph{Schriftsteller, Mediziner}!Reigen. Zehn Dialoge@\strich\emph{Reigen. Zehn Dialoge}|pwv} beendet
               hatte.\pend
           
\pstart
           Es hielt aber auch die Probe aus, wo dieſes Milieu und die entſprechende Stimmung
               fehlte: ich las es nämlich 2 Tage ſpäter in Alland\oindex{Alland@\textbf{Alland}, \emph{Verwaltungsgebiet}|pw} meinem Freunde Director von \textsc{Weismeyr\pwindex{Weismayr, Alexander von 24.\,9.\,1867 Wien – 10.\,3.\,1907 ebd.@\textsc{Weismayr, Alexander von} (24.\,9.\,1867 Wien – 10.\,3.\,1907 ebd.), \emph{Internist}|pw}} vor und fand de{\geminationn} erſten Eindruck vollauf beſtätigt
              und erzielte damit den vollſten Erfolg. Welches Paar aus dem Reigen\pwindex{Schnitzler, Arthur 15.\,5.\,1862 Wien – 21.\,10.\,1931 ebd.@\textsc{Schnitzler, Arthur} (15.\,5.\,1862 Wien – 21.\,10.\,1931 ebd.), \emph{Schriftsteller, Mediziner}!Reigen. Zehn Dialoge@\strich\emph{Reigen. Zehn Dialoge}|pw} mir am beſten gefällt, weiß ich eigentlich {\pb}nicht, vielleicht »der Gatte und das ſüße Mädl\pwindex{Schnitzler, Arthur 15.\,5.\,1862 Wien – 21.\,10.\,1931 ebd.@\textsc{Schnitzler, Arthur} (15.\,5.\,1862 Wien – 21.\,10.\,1931 ebd.), \emph{Schriftsteller, Mediziner}!Reigen. Zehn Dialoge@\strich\emph{Reigen. Zehn Dialoge}|pwv}«.\pend
           
\pstart
           Daß Sie das Buch\pwindex{Schnitzler, Arthur 15.\,5.\,1862 Wien – 21.\,10.\,1931 ebd.@\textsc{Schnitzler, Arthur} (15.\,5.\,1862 Wien – 21.\,10.\,1931 ebd.), \emph{Schriftsteller, Mediziner}!Reigen. Zehn Dialoge@\strich\emph{Reigen. Zehn Dialoge}|pwv} – ich will
               hoffen, nur \uline{vorläufig} – mit Ausſchluß der
               Öffentlichkeit erſcheinen ließen, thut mir ſehr leid, in Anbetracht deſſen, daß der
                  Verlag\orgindex{Albert Langen@Albert Langen|pwv} von \textsc{\uline{Langen}\pwindex{Langen, Albert 8.\,7.\,1869 Antwerpen – 30.\,4.\,1909 München@\textsc{Langen, Albert} (8.\,7.\,1869 Antwerpen – 30.\,4.\,1909 München), \emph{Verleger}|pw}} die ſtärkſten franz\oindex{Frankreich@\textbf{Frankreich}|pw}. Sachen bringt, glaube
               ich gewiß, daß er Ihren reno{\geminationm}irten Namen mit Vergnügen
               seinem Verlage\orgindex{Albert Langen@Albert Langen|pwv} einverleibt
               hätte. Ein \textcolor{gray}{Bomben}erfolg im Genre der »\label{K_L04227-4v}\edtext{\textsc{Lettres des femmes\pwindex{Prévost, Marcel 1.\,5.\,1862 Paris – 8.\,4.\,1941 Vianne@\textsc{Prévost, Marcel} (1.\,5.\,1862 Paris – 8.\,4.\,1941 Vianne), \emph{Schriftsteller}!Lettres de femmes@\strich\emph{Lettres de femmes}|pw}}}{\lemma{\textnormal{\emph{Lettres des femmes}}}\Cendnote{\textnormal{Die deutsche Übersetzung von Marcel Prévosts\pwindex{Prévost, Marcel 1.\,5.\,1862 Paris – 8.\,4.\,1941 Vianne@\textsc{Prévost, Marcel} (1.\,5.\,1862 Paris – 8.\,4.\,1941 Vianne), \emph{Schriftsteller}|pwk} Titel \emph{Lettres de Femmes}\pwindex{Prévost, Marcel 1.\,5.\,1862 Paris – 8.\,4.\,1941 Vianne@\textsc{Prévost, Marcel} (1.\,5.\,1862 Paris – 8.\,4.\,1941 Vianne), \emph{Schriftsteller}!Lettres de femmes@\strich\emph{Lettres de femmes}|pwk}, einer Sammlung fiktiver Briefe von
                  Frauen mit zum Teil libertinem und erotischen Inhalt, war 1895 als
                  eines der ersten Bücher im Verlagsprogramm von \emph{Albert Langen}\orgindex{Albert Langen@Albert Langen|pwk} erschienen (Marcel Prévost\pwindex{Prévost, Marcel 1.\,5.\,1862 Paris – 8.\,4.\,1941 Vianne@\textsc{Prévost, Marcel} (1.\,5.\,1862 Paris – 8.\,4.\,1941 Vianne), \emph{Schriftsteller}|pwk}: \emph{Pariserinnen. Autorisierte Übersetzung aus dem Französischen von
                        A. L.}\pwindex{Prévost, Marcel 1.\,5.\,1862 Paris – 8.\,4.\,1941 Vianne@\textsc{Prévost, Marcel} (1.\,5.\,1862 Paris – 8.\,4.\,1941 Vianne), \emph{Schriftsteller}!Pariserinnen@\strich\emph{Pariserinnen}|pwk}, Paris und Leipzig: \emph{Albert Langen}\orgindex{Albert Langen@Albert Langen|pwk}{ }1895).}}}\label{K_L04227-4}« ſcheint mir im Falle der Veröffentlichung ſicher geweſen zu{ }ſein.\pend
           
\pstart
           Ich weiß daher auch nicht, ob ich des Werkes\pwindex{Schnitzler, Arthur 15.\,5.\,1862 Wien – 21.\,10.\,1931 ebd.@\textsc{Schnitzler, Arthur} (15.\,5.\,1862 Wien – 21.\,10.\,1931 ebd.), \emph{Schriftsteller, Mediziner}!Reigen. Zehn Dialoge@\strich\emph{Reigen. Zehn Dialoge}|pwv} Erwähnung thun darf; ich habe nämlich eben einen
               Aufſatz »\label{K_L04227-5v}\edtext{\textsc{\uline{Das litterarische Wien\pwindex{Sosnosky, Theodor von 4.\,1.\,1866 Budapest – 3.\,2.\,1943 Wien@\textsc{Sosnosky, Theodor von} (4.\,1.\,1866 Budapest – 3.\,2.\,1943 Wien), \emph{Schriftsteller}!Jung-Wien«. Studienblätter@\strich\emph{»Jung-Wien«. Studienblätter}|pw}}}}{\lemma{\textnormal{\emph{Das litterarische Wien}}}\Cendnote{\textnormal{\emph{»Jung-Wien«. Studienblätter von Theodor v. Sosnosky}\pwindex{Sosnosky, Theodor von 4.\,1.\,1866 Budapest – 3.\,2.\,1943 Wien@\textsc{Sosnosky, Theodor von} (4.\,1.\,1866 Budapest – 3.\,2.\,1943 Wien), \emph{Schriftsteller}!Jung-Wien«. Studienblätter@\strich\emph{»Jung-Wien«. Studienblätter}|pwk}. In:
                     Beilage zur \emph{Norddeutschen Allgemeinen
                        Zeitung}\pwindex{Norddeutsche Allgemeine Zeitung@\emph{Norddeutsche Allgemeine Zeitung}|pwk}, Jg. 40, Nr. 17a, 20. 1. 1901,
                     S. 1–2.}}}\label{K_L04227-5}« für die »Norddeutſche allgemeine Ztg\pwindex{Norddeutsche Allgemeine Zeitung@\emph{Norddeutsche Allgemeine Zeitung}|pw}« unter der Feder, worin Sie, Herr Doktor,
               natürlich eine Rolle ſpielen; es wäre mir natürlich ſehr lieb, we{\geminationn} ich dieſes Ihr Meiſterwerk\pwindex{Schnitzler, Arthur 15.\,5.\,1862 Wien – 21.\,10.\,1931 ebd.@\textsc{Schnitzler, Arthur} (15.\,5.\,1862 Wien – 21.\,10.\,1931 ebd.), \emph{Schriftsteller, Mediziner}!Reigen. Zehn Dialoge@\strich\emph{Reigen. Zehn Dialoge}|pwv} wenigſtens flüchtig ſtreifen dürfte; ohne Ihre
              {\pb}ausdrückliche Einwilligung möchte ich’s aber nicht thun.\pend
           
\pstart
           Schluſſe will ich nur noch bemerken, daß ich bezüglich der Anthologie\pwindex{deutsche Lyrik des 19. Jahrhunderts. Eine poetische Revue@\emph{Die deutsche Lyrik des 19. Jahrhunderts. Eine poetische Revue}|pwv}, zu der Sie mir freundlicher
               Weiſe einige Gedichte\pwindex{Schnitzler, Arthur 15.\,5.\,1862 Wien – 21.\,10.\,1931 ebd.@\textsc{Schnitzler, Arthur} (15.\,5.\,1862 Wien – 21.\,10.\,1931 ebd.), \emph{Schriftsteller, Mediziner}!Wie wir so still@\strich\emph{Wie wir so still}|pwv}\pwindex{Schnitzler, Arthur 15.\,5.\,1862 Wien – 21.\,10.\,1931 ebd.@\textsc{Schnitzler, Arthur} (15.\,5.\,1862 Wien – 21.\,10.\,1931 ebd.), \emph{Schriftsteller, Mediziner}!Anfang vom Ende@\strich\emph{Anfang vom Ende}|pwv}\pwindex{Schnitzler, Arthur 15.\,5.\,1862 Wien – 21.\,10.\,1931 ebd.@\textsc{Schnitzler, Arthur} (15.\,5.\,1862 Wien – 21.\,10.\,1931 ebd.), \emph{Schriftsteller, Mediziner}!Ohnmacht@\strich\emph{Ohnmacht}|pwv} zur Verfügung geſtellt, noch nichts Definitives weiß; am 1. Mai gieng ſie an \textsc{Cotta}\orgindex{J.G. Cotta’sche Buchhandlung Nachfolger@J.G. Cotta’sche Buchhandlung Nachfolger|pw} ab, der ſich für
               dieſelbe intereßirte; seither sind \uline{erst} 1 ½ Monat
               vergangen; Verleger aber haben es in der Regel nicht eilig.\pend
           
\pstart
           Nochmals, geehrter Herr Doktor, meinen aufrichtigen Dank für die wirkliche
               Freude, die Sie mir durch das Buch\pwindex{Schnitzler, Arthur 15.\,5.\,1862 Wien – 21.\,10.\,1931 ebd.@\textsc{Schnitzler, Arthur} (15.\,5.\,1862 Wien – 21.\,10.\,1931 ebd.), \emph{Schriftsteller, Mediziner}!Reigen. Zehn Dialoge@\strich\emph{Reigen. Zehn Dialoge}|pwv} und die \label{K_L04227-6v}\edtext{liebenswürdige Widmung}{\lemma{\textnormal{\emph{liebenswürdige Widmung}}}\Cendnote{\textnormal{nicht überliefert}}}\label{K_L04227-6} bereitet haben{\\[\baselineskip]} Ihrem Sie hochſchätzenden{\\[\baselineskip]}\spacefill\mbox{Theodor vSosnosky}\pend
           \leftskip=0em{}\selectlanguage{ngerman}\endnumbering\briefempfaengerindex{Schnitzler, Arthur@\textsc{Schnitzler, Arthur}!zzzSosnosky, Theodor von@\emph{von Theodor von Sosnosky}!1900-06-221@{22. 6. 1900}|)be}\mylabel{L04227h}
\begin{anhang}
\end{anhang}\newcommand{\dateiname}{L04227}\newcommand{\titel}{Theodor von Sosnosky an Arthur Schnitzler, 22. 6. 1900}\newcommand{\editorInnen}{Herausgegeben von Jahnke, SelmaMüller, Martin Anton}\input{../tex-inputs/latex-leseansicht-abspann}
